\documentclass[aspectratio=169,10pt]{beamer}



% ============================================================

% PACKAGES

% ============================================================



\usepackage[utf8]{inputenc}

\usepackage[T1]{fontenc}

\usepackage[brazil]{babel}

\usepackage{lmodern}



\usepackage{graphicx}

\usepackage{booktabs}

\usepackage{tabularx}

\usepackage{array}

\usepackage{tikz}

\usepackage{amsmath,amssymb}

\usepackage{ragged2e}

\usepackage{microtype}

\usepackage{multirow}



\usetikzlibrary{

&#x20;   calc,

&#x20;   arrows.meta,

&#x20;   positioning,

&#x20;   shapes.geometric

}



\setbeamertemplate{navigation symbols}{}



\setbeamersize{

&#x20;   text margin left=0.72cm,

&#x20;   text margin right=0.72cm

}



% ============================================================

% ACADEMIC INFORMATION

% ============================================================



\newcommand{\graduate}{

Department of Industrial Engineering

}



\newcommand{\course}{

IND2097 - SPECIAL TOP IN OP RESEARCH - 2026.2

}



\newcommand{\theme}{

Homework 1 --- Decision Rules

}



\newcommand{\student}{

Carlos Souza

}



\newcommand{\registration}{

2612398

}



\newcommand{\semester}{

3RA --- 2026.2

}



\newcommand{\prof}{

BRUNO FANZERES DOS SANTOS

}



\newcommand{\puclogo}{

puc_logo.png

}



\newcommand{\seminartitle}{

Regularized Linear Decision Rules

}



\newcommand{\seminarsubtitle}{

A Data-Driven Application to Residential Battery Control

}



\hypersetup{

&#x20;   pdftitle={

&#x20;       \course\space --- Regularized Linear Decision Rules

&#x20;   },

&#x20;   pdfauthor={

&#x20;       \student

&#x20;   }

}



% ============================================================

% COLORS

% ============================================================



\definecolor{mainblue}{RGB}{18,55,135}

\definecolor{darkblue}{RGB}{8,38,100}

\definecolor{accentgold}{RGB}{205,158,10}



\definecolor{lightgray}{RGB}{242,244,248}

\definecolor{textgray}{RGB}{40,45,55}



\definecolor{goodgreen}{RGB}{19,111,83}



\definecolor{softblue}{RGB}{235,241,250}

\definecolor{softgreen}{RGB}{231,245,239}

\definecolor{softgold}{RGB}{250,246,225}



% ============================================================

% BEAMER STYLE

% ============================================================



\setbeamercolor{normal text}{

&#x20;   fg=textgray,

&#x20;   bg=white

}



\setbeamercolor{frametitle}{

&#x20;   fg=mainblue,

&#x20;   bg=white

}



\setbeamercolor{framesubtitle}{

&#x20;   fg=darkblue,

&#x20;   bg=white

}



\setbeamerfont{frametitle}{

&#x20;   series=\bfseries,

&#x20;   size=\Large

}



\setbeamerfont{framesubtitle}{

&#x20;   size=\small

}



\setbeamerfont{footline}{

&#x20;   size=\scriptsize

}



\setbeamercolor{block title}{

&#x20;   bg=lightgray,

&#x20;   fg=mainblue

}



\setbeamercolor{block body}{

&#x20;   bg=lightgray,

&#x20;   fg=textgray

}



% ============================================================

% FRAME TITLE

% ============================================================



\setbeamertemplate{frametitle}{%



\begin{tikzpicture}[remember picture,overlay]



&#x20;   \fill[mainblue]

&#x20;   (current page.north west)

&#x20;   rectangle ++(0.45,-1.10);



&#x20;   \node[

&#x20;       white,

&#x20;       font=\bfseries\tiny

&#x20;   ]

&#x20;   at ($(current page.north west)+(0.225,-0.55)$)

&#x20;   {STOR};



\end{tikzpicture}%



\vspace{0.12cm}



\hspace{0.05cm}

\insertframetitle

\par



\vspace{-0.05cm}



{

&#x20;   \usebeamerfont{framesubtitle}

&#x20;   \usebeamercolor[fg]{framesubtitle}



&#x20;   \hspace{0.05cm}

&#x20;   \insertframesubtitle

&#x20;   \par

}



\vspace{0.15cm}

}



% ============================================================

% FOOTLINE

% ============================================================



\setbeamertemplate{footline}{%



\leavevmode%



\hbox{%



\begin{beamercolorbox}[

&#x20;   wd=.72\paperwidth,

&#x20;   ht=2.5ex,

&#x20;   dp=1ex,

&#x20;   leftskip=0.7cm

]{footline}

\end{beamercolorbox}%



\begin{beamercolorbox}[

&#x20;   wd=.28\paperwidth,

&#x20;   ht=2.5ex,

&#x20;   dp=1ex,

&#x20;   rightskip=0.7cm plus1fil

]{footline}



\hfill

\insertframenumber/\inserttotalframenumber



\end{beamercolorbox}



}

}



% ============================================================

% CUSTOM BLOCK

% ============================================================



\newenvironment{blueblock}[1]{%



\begin{beamercolorbox}[

&#x20;   rounded=true,

&#x20;   shadow=false,

&#x20;   sep=0.15cm,

&#x20;   wd=\linewidth

]{block title}



\color{mainblue}

\bfseries

\#1



\end{beamercolorbox}



\vspace{-0.04cm}



\begin{beamercolorbox}[

&#x20;   rounded=true,

&#x20;   shadow=false,

&#x20;   sep=0.18cm,

&#x20;   wd=\linewidth,

&#x20;   vmode

]{block body}



\small



}{

\end{beamercolorbox}

}



% ============================================================

% OUTPUT PLACEHOLDER

% If output exists -> show it.

% Otherwise -> presentation still compiles.

% ============================================================



\newcommand{\resultfig}[2][0.96]{%



\IfFileExists{graficos/apresentacao/png/#2}

{

&#x20;   \includegraphics[

&#x20;       width=#1\linewidth

&#x20;   ]{graficos/apresentacao/png/#2}

}

{

&#x20;   \fbox{

&#x20;       \parbox[c][0.30\textheight][c]{0.90\linewidth}{

&#x20;           \centering

&#x20;           \small



&#x20;           Output ainda não gerado.



&#x20;           \vspace{0.10cm}



&#x20;           \texttt{graficos/apresentacao/png/#2}

&#x20;       }

&#x20;   }

}



}



% ============================================================

% DOCUMENT

% ============================================================



\begin{document}



% ============================================================

% CAPA

% ============================================================



{

\setbeamertemplate{footline}{}



\begin{frame}[plain]



\begin{tikzpicture}[remember picture,overlay]



% Fundo

\fill[white]

(current page.south west)

rectangle

(current page.north east);



% Área diagonal azul

\fill[mainblue]

($(current page.north east)+(-4.60,0)$)

\--

(current page.north east)

\--

(current page.south east)

\--

($(current page.south east)+(-2.10,0)$)

\-- cycle;



% Bloco STOR

\fill[mainblue]

($(current page.north west)+(0.55,-0.35)$)

rectangle ++(0.80,-1.25);



\node[

&#x20;   white,

&#x20;   font=\bfseries\large

]

at ($(current page.north west)+(0.95,-0.98)$)

{STOR};



% Logo

\node[

&#x20;   anchor=north east,

&#x20;   inner sep=0pt

]

at ($(current page.north east)+(-0.50,-0.42)$)

{

&#x20;   \includegraphics[

&#x20;       height=2.25cm

&#x20;   ]{\puclogo}

};



% Lateral

\node[

&#x20;   white,

&#x20;   align=center,

&#x20;   font=\bfseries\large

]

at ($(current page.south east)+(-1.45,1.18)$)

{

&#x20;   Decision\\\\

&#x20;   Rules

};



% Conteúdo

\node[

&#x20;   anchor=north west,

&#x20;   align=left,

&#x20;   text width=10.15cm

]

at ($(current page.north west)+(0.65,-1.62)$)

{%



{\color{mainblue}

\bfseries

\fontsize{21}{24}\selectfont

Regularized Linear Decision Rules

}

\\\\[0.08cm]



{\color{darkblue}

\bfseries

\fontsize{13.5}{16}\selectfont

Multistage Stochastic Programming

}

\\\\[0.16cm]



{\color{textgray}

\small

Based on Nazare \& Street (2023)

}

\\\\[0.28cm]



{\color{mainblue}

\bfseries

\normalsize

Aplicação prática:

controle de bateria residencial

}

\\\\[0.42cm]



{\small

\textbf{Student:}

\student

\qquad

\textbf{Registration:}

\registration

}

\\\\[0.12cm]



{\small

\textbf{Professor:}

\prof

}

\\\\[0.12cm]



{\small

\textbf{Department:}

\graduate

}

\\\\[0.12cm]



{\small

\textbf{Course:}

\course

}

\\\\[0.12cm]



{\small

\textbf{Class / Semester:}

\semester

}



};



\end{tikzpicture}



\end{frame}



}



% ============================================================

% SLIDE 1 -- HISTÓRIA

% ============================================================



\begin{frame}

{O problema}

{uma decisão simples que depende de um futuro incerto !!}



\begin{center}



\includegraphics[

&#x20;   width=0.88\linewidth,

&#x20;   height=0.55\textheight,

&#x20;   keepaspectratio

]{casa.png}



\end{center}



\vspace{-0.12cm}



\begin{blueblock}{A situação}



A cada meia hora, precisamos decidir:



\\[

\boxed{

\text{usar energia}

\quad

\text{armazenar}

\quad

\text{ou}

\quad

\text{comprar da rede}

}

\\]



Mas ainda não conhecemos exatamente o consumo

e a geração solar das próximas horas.



\end{blueblock}



\end{frame}



% ============================================================

% SLIDE 2 -- MOTIVAÇÃO

% ============================================================



\begin{frame}

{Por que esse problema é difícil?}

{uma decisão agora altera as opções disponíveis depois}



\small

\abovedisplayskip=4pt

\belowdisplayskip=4pt



\begin{columns}[T,totalwidth=\textwidth]



\begin{column}{0.47\textwidth}



\begin{blueblock}{Hoje}



Às 10h observamos:



\\[

GC\_{10},

\qquad
GG\_{10},

\qquad

B\_{10}.

\\]



Precisamos decidir quanto:



\\[

\text{carregar},

\qquad

\text{descarregar}.

\\]



\end{blueblock}



\end{column}



\begin{column}{0.49\textwidth}



\begin{blueblock}{Amanhã... ou daqui a algumas horas}



A decisão atual modifica:



\\[

B\_{11},

B\_{12},

\ldots

\\]



e portanto altera as decisões

que estarão disponíveis no futuro.Além disso,



\\[

GC\_{t+1},

GG\_{t+1}

\\]



ainda são incertos.



\end{blueblock}



\end{column}



\end{columns}



\vspace{0.14cm}



\begin{beamercolorbox}[

&#x20;   rounded=true,

&#x20;   sep=0.16cm

]{block body}



\centering



\large

\color{mainblue}

\bfseries



Como tomar boas decisões hoje

sem conhecer exatamente o futuro?



\end{beamercolorbox}



\end{frame}



% ============================================================

% SLIDE 3 -- PONTE PARA TEORIA

% ============================================================



\begin{frame}

{Da história para a matemática}

{o problema pode ser interpretado como um problema de programação matemática multiestágio}



\begin{center}



\begin{tikzpicture}[

&#x20;   \>=Latex,

&#x20;   node distance=0.55cm,

&#x20;   every node/.style={

&#x20;       align=center,

&#x20;       font=\small

&#x20;   }

]



\node[

&#x20;   draw=mainblue,

&#x20;   fill=softblue,

&#x20;   rounded corners,

&#x20;   minimum width=2.0cm,

&#x20;   minimum height=0.8cm

] (x1)

{

Decisão\\\\

$x_1$

};



\node[

&#x20;   draw=accentgold,

&#x20;   fill=softgold,

&#x20;   circle,

&#x20;   right=of x1,

&#x20;   minimum size=1.05cm

] (xi2)

{

$\xi_2$

};



\node[

&#x20;   draw=mainblue,

&#x20;   fill=softblue,

&#x20;   rounded corners,

&#x20;   right=of xi2,

&#x20;   minimum width=2.0cm,

&#x20;   minimum height=0.8cm

] (x2)

{

Decisão\\\\

$x_2$

};



\node[

&#x20;   draw=accentgold,

&#x20;   fill=softgold,

&#x20;   circle,

&#x20;   right=of x2,

&#x20;   minimum size=1.05cm

] (xi3)

{

$\xi_3$

};



\node[

&#x20;   draw=mainblue,

&#x20;   fill=softblue,

&#x20;   rounded corners,

&#x20;   right=of xi3,

&#x20;   minimum width=2.0cm,

&#x20;   minimum height=0.8cm

] (x3)

{

Decisão\\\\

$x_3$

};



\node[

&#x20;   right=0.55cm of x3

]

{

$\cdots$

};



\draw[->,thick] (x1) -- (xi2);

\draw[->,thick] (xi2) -- (x2);

\draw[->,thick] (x2) -- (xi3);

\draw[->,thick] (xi3) -- (x3);



\end{tikzpicture}



\end{center}



\vspace{0.25cm}



\begin{blueblock}{Política de decisão}



No instante $t$, usamos apenas a informação

que já foi observada:



\\[

\boxed{

x_t

\=

\pi_t(

\xi_1,\ldots,\xi_t

)

}

\\]



Essa condição é chamada de

\textbf{não antecipatividade}.



\end{blueblock}



\end{frame}



% ============================================================

% SLIDE 4 -- TEORIA DO PAPER

% ============================================================



\begin{frame}[shrink=12]

{Base teórica do artigo}

{do problema multiestágio à regra de decisão linear}



\small

\abovedisplayskip=4pt

\belowdisplayskip=4pt



\begin{blueblock}{Problema}



Queremos tomar decisões sequenciais sob incerteza:



\\[

\boxed{

x_t = \pi_t(\xi_1,\ldots,\xi_t)

}

\\]



onde $x_t$ é a decisão no estágio $t$ e

$\xi_1,\ldots,\xi_t$ representam as informações observadas até $t$.



\end{blueblock}



\vspace{0.06cm}



\begin{blueblock}{Dificuldade}



Encontrar diretamente a política $\pi_t(\cdot)$ é difícil.



Então o artigo aproxima essa política por uma

\textbf{Linear Decision Rule (LDR)}:



\\[

\boxed{

x_t =

\sum_k

\phi_k(\xi\_{1:t})\\,\theta\_{t,k}

}

\\]



A ideia é substituir uma política geral por uma regra parametrizada

pelos coeficientes $\theta$.



\end{blueblock}



\vspace{0.06cm}



\centering

\small

A vantagem é que o problema fica mais tratável computacionalmente.



\end{frame}



% ============================================================

% SLIDE 5 -- ADALASSO

% ============================================================



\begin{frame}

{Problema central do artigo}

{overfitting e regularização via AdaLASSO}



\small

\abovedisplayskip=3pt

\belowdisplayskip=3pt



\begin{columns}[T,totalwidth=\textwidth]



\begin{column}{0.44\textwidth}



\begin{blueblock}{Risco de overfitting}



Se o número de coeficientes da LDR for grande em relação ao número de cenários de treino,



\\[

\boxed{

p = \\#\theta \approx N

}

\\]



então a política pode ajustar muito bem os dados de treino,

mas apresentar baixo desempenho fora da amostra.



\\[

\boxed{

p \approx N

\\;\Longrightarrow\\;

\text{overfitting}

}

\\]



\end{blueblock}



\end{column}



\begin{column}{0.52\textwidth}



\begin{blueblock}{Solução proposta no artigo}



\footnotesize



O artigo adiciona regularização AdaLASSO:



\\[

\boxed{

\min\_{\theta}

\\; C\_{\mathrm{train}}(\theta)

\+

\lambda

\sum_j

\frac{|\theta_j|}{|\theta_j^{(0)}|}

}

\\]



onde:



\begin{itemize}\itemsep1pt\parsep0pt\topsep2pt

&#x20;   \item $C\_{\mathrm{train}}(\theta)$: custo no treinamento;

&#x20;   \item $\lambda$: intensidade da penalização;

&#x20;   \item $\theta^{(0)}$: coeficientes da LDR sem regularização.

\end{itemize}



A meta é obter uma política mais simples e com melhor desempenho OOS.



\vspace{0.05cm}

\centering

$

\boxed{

\text{LDR simples}

\\;\Rightarrow\\;

\text{melhor generalização?}

}

$



\end{blueblock}



\end{column}



\end{columns}



\end{frame}



% ============================================================

% PONTE -- PAPER E EXPERIMENTO

% ============================================================



\begin{frame}

{Do paper X nosso experimento}

{a mesma lógica em um problema diferente}



\small

\abovedisplayskip=3pt

\belowdisplayskip=3pt



\begin{columns}[T,totalwidth=\textwidth]



\begin{column}{0.51\textwidth}



\begin{blueblock}{Equivalência entre os problemas}



\centering

\scriptsize

\renewcommand{\arraystretch}{1.18}



\begin{tabular}{@{}lcl@{}}

\toprule

\textbf{Paper} & & \textbf{Ausgrid} \\\\

\midrule

Reservatório & $\leftrightarrow$ & Bateria \\\\

Água armazenada & $\leftrightarrow$ & $B_t$ \\\\

Afluência & $\leftrightarrow$ & $GC_t-GG_t$ \\\\

Turbinamento & $\leftrightarrow$ & descarga \\\\

Geração térmica & $\leftrightarrow$ & compra da rede \\\\

Custo térmico & $\leftrightarrow$ & conta de energia \\\\

\bottomrule

\end{tabular}



\end{blueblock}



\end{column}



\begin{column}{0.45\textwidth}



\begin{blueblock}{Como usamos a LDR aqui}



\footnotesize

No paper, a LDR define um estado-alvo para o reservatório.



No nosso caso, ela define um estado-alvo para a bateria:



\\[

\boxed{

B_t^{\text{target}}

\=

\theta\_{t0}

\+

\sum\_{\ell,k}

\theta\_{t\ell k}\\,z\_{t-\ell}^{\\,k}

}

\\]



onde $z\_{t-\ell}$ representa informações recentes da carga líquida.



\end{blueblock}



\end{column}



\end{columns}



\vspace{0.12cm}



\begin{beamercolorbox}[

&#x20;   rounded=true,

&#x20;   sep=0.18cm,

&#x20;   wd=\linewidth

]{block body}



\centering

\small

\textbf{Em português:}\quad

\emph{“Olhando o consumo e a geração solar recentes,

quanto de energia eu gostaria de ter guardado na bateria agora?”}



\end{beamercolorbox}



\end{frame}



% ============================================================

% SLIDE 6 -- METODOLOGIA

% ============================================================



\begin{frame}

{Nossa metodologia}

{dados reais do Ausgrid + bateria simulada}



\begin{center}



\begin{tikzpicture}[

&#x20;   \>=Latex,

&#x20;   node distance=0.28cm,

&#x20;   every node/.style={

&#x20;       align=center,

&#x20;       font=\scriptsize

&#x20;   }

]



\node[

&#x20;   draw=mainblue,

&#x20;   fill=softblue,

&#x20;   rounded corners,

&#x20;   minimum width=1.55cm,

&#x20;   minimum height=0.80cm

] (data)

{

Ausgrid\\\\

$GC_t,GG_t$

};



\node[

&#x20;   draw=mainblue,

&#x20;   fill=softblue,

&#x20;   rounded corners,

&#x20;   right=of data,

&#x20;   minimum width=1.55cm,

&#x20;   minimum height=0.80cm

] (net)

{

Carga líquida\\\\

$L_t=GC_t-GG_t$

};



\node[

&#x20;   draw=mainblue,

&#x20;   fill=softblue,

&#x20;   rounded corners,

&#x20;   right=of net,

&#x20;   minimum width=1.55cm,

&#x20;   minimum height=0.80cm

] (train)

{

Treino\\\\

LDR

};



\node[

&#x20;   draw=goodgreen,

&#x20;   fill=softgreen,

&#x20;   rounded corners,

&#x20;   right=of train,

&#x20;   minimum width=1.65cm,

&#x20;   minimum height=0.80cm

] (ada)

{

AdaLASSO\\\\

vários $\lambda$

};



\node[

&#x20;   draw=accentgold,

&#x20;   fill=softgold,

&#x20;   rounded corners,

&#x20;   right=of ada,

&#x20;   minimum width=1.55cm,

&#x20;   minimum height=0.80cm

] (val)

{

Validação\\\\

$\lambda^\star$

};



\node[

&#x20;   draw=goodgreen,

&#x20;   fill=softgreen,
&#x20;   rounded corners,

&#x20;   right=of val,

&#x20;   minimum width=1.55cm,

&#x20;   minimum height=0.80cm

] (test)

{

Teste\\\\

OOS

};



\draw[->,thick] (data) -- (net);

\draw[->,thick] (net) -- (train);

\draw[->,thick] (train) -- (ada);

\draw[->,thick] (ada) -- (val);

\draw[->,thick] (val) -- (test);



\end{tikzpicture}



\end{center}



\vspace{0.32cm}



\begin{columns}[T,totalwidth=\textwidth]



\begin{column}{0.47\textwidth}



\begin{blueblock}{Dados}



\\[

GC_t

\=

\text{consumo}

\\]



\\[

GG_t

\=

\text{geração solar}

\\]



\\[

L_t

\=

GC_t-GG_t.

\\]



Cada dia possui:



\\[

T=48

\\]



períodos de 30 minutos.



\end{blueblock}



\end{column}



\begin{column}{0.49\textwidth}



\begin{blueblock}{Avaliação}



Escolhemos:



\\[

\boxed{

\lambda^\star

\=

\arg\min\_\lambda

C\_{\mathrm{val}}(\lambda)

}

\\]



e só depois comparamos

as políticas no teste OOS.



\end{blueblock}



\end{column}



\end{columns}



\end{frame}



% ============================================================

% SLIDE 7 -- MODELO BATERIA

% ============================================================



\begin{frame}

{Modelo da bateria}

{estado, controle e função objetivo}



\begin{columns}[T,totalwidth=\textwidth]



\begin{column}{0.48\textwidth}



\begin{blueblock}{Estado}



A energia armazenada evolui como:



\\[

\boxed{

B_t

\=

B\_{t-1}

\+

\eta_c c_t

\-

\frac{d_t}{\eta_d}

}

\\]



onde:



\\[

c_t=\text{carga},

\qquad

d_t=\text{descarga}.

\\]



Com:



\\[

0\leq B_t\leq B\_{\max}.

\\]



\end{blueblock}



\end{column}



\begin{column}{0.48\textwidth}



\begin{blueblock}{Objetivo}



A energia comprada da rede é:



\\[

G_t.

\\]



Queremos minimizar:



\\[

\boxed{

\min

\mathbb E

\left[

\sum\_{t=1}^{T}

p_tG_t

\right]

}

\\]



respeitando todas as restrições

físicas da bateria.



\end{blueblock}



\end{column}



\end{columns}



\end{frame}



% ============================================================

% SLIDE 8 -- IMPLEMENTAÇÃO

% ============================================================



\begin{frame}

{Implementação}

{Julia, JuMP e Gurobi}



\begin{columns}[T,totalwidth=\textwidth]



\begin{column}{0.42\textwidth}



\begin{blueblock}{Tecnologias}



\begin{center}



{\Large

\textbf{Julia}

}



\\[

\Downarrow

\\]



{\Large

\textbf{JuMP}

}



\\[

\Downarrow

\\]



{\Large

\textbf{Gurobi}

}



\end{center}



\end{blueblock}



\end{column}



\begin{column}{0.54\textwidth}



\begin{blueblock}{Rotina experimental}



\begin{enumerate}



&#x20;   \item Estimar LDR sem regularização;



&#x20;   \item calcular pesos AdaLASSO;



&#x20;   \item resolver o modelo para vários $\lambda$;



&#x20;   \item escolher $\lambda^\star$ na validação;



&#x20;   \item avaliar LDR e AdaLASSO

&#x20;   na mesma amostra OOS.



\end{enumerate}



\end{blueblock}



\end{column}



\end{columns}



\end{frame}



% ============================================================

% SLIDE 9 -- RESULTADO 1

% ============================================================



\begin{frame}

{Resultados}

{o que acontece quando aumentamos a regularização?}



\begin{columns}[T,totalwidth=\textwidth]



\begin{column}{0.49\textwidth}



\centering



\resultfig{fig_06_treino_vs_validacao.png}



\vspace{0.05cm}



\scriptsize

Custo no treino e validação para diferentes valores de $\lambda$.



\end{column}



\begin{column}{0.49\textwidth}



\centering



\resultfig{fig_03_parcimonia_coeficientes.png}



\vspace{0.05cm}



\scriptsize

Número de coeficientes não nulos da política.



\end{column}



\end{columns}



\vspace{0.15cm}



\begin{blueblock}{Pergunta}



A redução da complexidade da política

também melhora seu desempenho fora da amostra?



\end{blueblock}



\end{frame}



% ============================================================

% SLIDE 10 -- RESULTADO OOS

% ============================================================



\begin{frame}

{Desempenho fora da amostra}

{LDR versus AdaLASSO-LDR}



\begin{columns}[T,totalwidth=\textwidth]



\begin{column}{0.49\textwidth}



\centering



\resultfig{fig_04_distribuicao_custos_oos.png}



\end{column}



\begin{column}{0.49\textwidth}



\centering



\resultfig{fig_07_custo_medio_oos.png}



\end{column}



\end{columns}



\vspace{0.15cm}



\begin{blueblock}{Interpretação}



O resultado principal não deve ser julgado

pelo custo obtido no treinamento,

mas pelo comportamento da política

em dias que não foram utilizados para estimá-la.



\end{blueblock}



\end{frame}



% ============================================================

% SLIDE 11 -- OUTPUT OPERACIONAL

% ============================================================



\begin{frame}

{O que a política realmente faz?}

{interpretação operacional da bateria}



\begin{columns}[T,totalwidth=\textwidth]



\begin{column}{0.59\textwidth}



\centering



\resultfig[0.98]{fig_05_trajetoria_bateria.png}



\end{column}



\begin{column}{0.37\textwidth}



\begin{blueblock}{Leitura}



A política determina

um estado-alvo:



\\[

B_t^{target}.

\\]



A bateria então:



\begin{itemize}



&#x20;   \item armazena energia em alguns períodos;



&#x20;   \item descarrega em outros;



&#x20;   \item reduz ou desloca compras da rede.



\end{itemize}



\end{blueblock}



\end{column}



\end{columns}



\end{frame}



% ============================================================

% SLIDE 12 -- PAPER BENCHMARK

% ============================================================



\begin{frame}

{Ligação com os resultados do artigo}

{o que Nazare \& Street encontraram}



\begin{columns}[T,totalwidth=\textwidth]



\begin{column}{0.47\textwidth}



\begin{blueblock}{Custo OOS}



No primeiro estudo de caso:



\\[

C\_{\mathrm{LDR}}

\=

2.00\text{ M\\$}

\\]



e, com regularização:



\\[

C\_{\mathrm{Ada}}

\=

1.63\text{ M\\$}.

\\]



Redução:



\\[

\boxed{

18.5\\%

}

\\]



\end{blueblock}



\end{column}



\begin{column}{0.49\textwidth}



\begin{blueblock}{Parcimônia}



Coeficientes não nulos:



\\[

45.0\\%

\\]



sem regularização,



\\[

\Downarrow

\\]



\\[

5.1\\%

\\]



na política escolhida.



Redução:



\\[

\boxed{

88.6\\%

}

\\]



\end{blueblock}



\end{column}



\end{columns}



\vspace{0.15cm}



\centering

\scriptsize



Fonte: Nazare \& Street (2023).



\end{frame}



% ============================================================

% SLIDE 13 -- FUTUROS PASSOS

% ============================================================



\begin{frame}

{Próximos passos}

{como expandir o experimento}



\begin{columns}[T,totalwidth=\textwidth]



\begin{column}{0.48\textwidth}



\begin{blueblock}{Dados}



\begin{itemize}



&#x20;   \item utilizar os três anos completos do Ausgrid;



&#x20;   \item repetir o experimento

&#x20;   para várias residências;



&#x20;   \item avaliar diferentes perfis

&#x20;   de consumo e geração solar.



\end{itemize}



\end{blueblock}



\end{column}



\begin{column}{0.48\textwidth}



\begin{blueblock}{Modelo}



\begin{itemize}



&#x20;   \item utilizar tarifas reais;



&#x20;   \item testar outras capacidades de bateria;



&#x20;   \item comparar com outros métodos,
&#x20;   como SDDP ou políticas alternativas.



\end{itemize}



\end{blueblock}



\end{column}



\end{columns}



\end{frame}



% ============================================================

% CONCLUSÃO

% ============================================================



\begin{frame}

{Mensagem final}

{decidir bem não significa apenas ajustar bem o treinamento}



\begin{center}



\vspace{0.35cm}



\\[

\boxed{

\text{Decisão sob incerteza}

}

\\]



\\[

\Downarrow

\\]



\\[

\boxed{

\text{Linear Decision Rules}

}

\\]



\\[

\Downarrow

\\]



\\[

\boxed{

\text{Regularização}

}

\\]



\\[

\Downarrow

\\]



\\[

\boxed{

\text{Generalização OOS}

}

\\]



\vspace{0.35cm}



\begin{beamercolorbox}[

&#x20;   rounded=true,

&#x20;   sep=0.30cm,

&#x20;   wd=0.86\textwidth

]{block body}



\centering



\Large

\color{mainblue}

\bfseries



Uma política que usa menos parâmetros

pode tomar decisões melhores

em situações que nunca observou?



\end{beamercolorbox}



\end{center}



\end{frame}



% ============================================================

% REFERENCES

% ============================================================



\begin{frame}

{Referências}

{}



\small



\begin{thebibliography}{9}



\bibitem{nazare2023}

F. Nazare and A. Street.



\newblock

\textit{Solving Multistage Stochastic Linear Programming via

Regularized Linear Decision Rules:

An Application to Hydrothermal Dispatch Planning}.



\newblock

European Journal of Operational Research,

309, 345--358, 2023.





\bibitem{bodur2018}

M. Bodur and J. Luedtke.



\newblock

\textit{Two-Stage Linear Decision Rules

for Multi-Stage Stochastic Programming}.





\bibitem{zou2006}

H. Zou.



\newblock

\textit{The Adaptive Lasso and Its Oracle Properties}.



\newblock

Journal of the American Statistical Association,

2006\.



\end{thebibliography}



\end{frame}



\end{document}